\section{Exact prime completions and the obstruction to deletion descent}
\label{app:finalprimeaudit}

This appendix studies the last prime of a solution. Once the other
prime factors are fixed, the quotient and the last prime are
determined (Theorem~\ref{thm:auditcriterion}), so Lehmer's conjecture
is equivalent to the assertion of Problem~\ref{prob:universalprime},
which we do not prove. We record what the exact divisibility already
implies, and show that deleting the largest prime of a solution never
gives a smaller solution.

\subsection{Notation and the unique possible last prime}
Let $S$ be a nonempty finite set of distinct odd primes and put
\begin{equation}\label{eq:auditprefix}
 A=\prod_{p\in S}p,\qquad B=\prod_{p\in S}(p-1)=\varphi(A),
 \qquad r=\max S.
\end{equation}
These are integers, $A$ is odd, and $B$ is positive and even. Moreover
$A/B>1$ and $A/B$ is not an integer: an even integer cannot divide the
odd integer $A$. For an integer $M\ge2$ write
\begin{equation}\label{eq:auditdefect}
 D=MB-A.
\end{equation}
An additional prime $q>r$ is different from every element of $S$.
Consequently its totient contribution is exactly $q-1$, and
\begin{align}
 \varphi(Aq)\mid Aq-1\text{ with quotient }M
 &\quad\Longleftrightarrow\quad Aq-1=MB(q-1)\notag\\
 &\quad\Longleftrightarrow\quad (MB-A)q=MB-1\notag\\
 &\quad\Longleftrightarrow\quad D(q-1)=A-1.
 \label{eq:auditidentity}
\end{align}
The last line forces $D>0$, since $A>1$ and $q>1$.
There has been no division by a possibly zero integer in obtaining this
conclusion.

\begin{lemma}[The quotient is forced]\label{lem:auditrounding}
If $q>r$ is prime and $Aq-1=MB(q-1)$, then
\begin{equation}\label{eq:auditrounding}
 0<M-\frac AB=\frac{A-1}{B(q-1)}<\frac12,
 \qquad M=\left\lceil\frac AB\right\rceil.
\end{equation}
\end{lemma}
\begin{proof}
Every factor $j/(j-1)$ is greater than $1$ for an integer $j\ge3$.
Since $S$ is a subset of $\{3,4,\ldots,r\}$, inserting its missing
integers can only increase the product. Hence
\[
 \frac AB=\prod_{p\in S}\frac p{p-1}
 \le\prod_{j=3}^{r}\frac j{j-1}=\frac r2.
\]
The two primes $q>r$ are odd, so $q\ge r+2$, and therefore
$r/(q-1)<1$. The identity $D(q-1)=A-1$ now gives
\[
 0<\frac DB=\frac{A-1}{B(q-1)}
 <\frac{A/B}{q-1}\le\frac{r}{2(q-1)}<\frac12.
\]
By $D=MB-A$ this is the stated estimate. An integer strictly between
$A/B$ and $A/B+1/2$ must be its ceiling.
\end{proof}

\begin{theorem}[Exact final-prime criterion]\label{thm:auditcriterion}
For $S$ as in \eqref{eq:auditprefix}, define
\[
 M=\left\lceil\frac AB\right\rceil,\qquad D=MB-A.
\]
There exists a Lehmer composite with prime support $S\cup\{q\}$ and
$q>\max S$ if and only if
\begin{equation}\label{eq:auditcertificate}
 D\mid A-1,\qquad
 q=1+\frac{A-1}{D}>r,\qquad q\text{ is prime}.
\end{equation}
The possible prime $q$ is unique.
\end{theorem}
\begin{proof}
The observations above give $M\ge2$ and $D>0$ for this choice of
$M$. If a Lehmer composite exists, Lemma~\ref{lem:auditrounding}
forces this $M$, and \eqref{eq:auditidentity} forces every condition
in \eqref{eq:auditcertificate}. Conversely, suppose those conditions
hold. The primality and ordering of $q$ imply
$\varphi(Aq)=B(q-1)$. The definition of $q$ gives
$Dq=D+A-1=MB-1$, and hence
\[
 MB(q-1)=(A+D)q-MB=Aq-1.
\]
Thus $\varphi(Aq)\mid Aq-1$. Since $S$ is nonempty and $q>r$,
$Aq$ has at least two distinct prime factors and is composite.
Uniqueness follows from the displayed formula for $q$.
\end{proof}

\begin{figure}[tbp]
\centering
\begin{tikzpicture}[>=Latex, every node/.style={font=\small},
 box/.style={draw=blue!45!black,rounded corners=2pt,fill=blue!3,
 text width=10.5cm,align=center,inner sep=6pt}]
 \node[box] (s) {Distinct odd primes $S$, with $A=\prod p$, $B=\prod(p-1)$};
 \node[box,below=6mm of s] (m) {$M=\lceil A/B\rceil$, $D=MB-A>0$};
 \node[box,below=6mm of m] (d) {Exact integer gate: $D\mid A-1$};
 \node[box,below=6mm of d] (q) {$q=1+(A-1)/D$: require $q>\max S$ and $q$ prime};
 \node[box,below=6mm of q] (n) {An accepted certificate gives $Aq-1=M\varphi(Aq)$};
 \draw[->] (s)--(m); \draw[->] (m)--(d);
 \draw[->] (d)--(q); \draw[->] (q)--(n);
\end{tikzpicture}
\caption{The completion criterion of Theorem~\ref{thm:auditcriterion}.
Lehmer's conjecture asserts that the last two tests never both pass
(Problem~\ref{prob:universalprime}).}
\label{fig:auditgates}
\end{figure}

\subsection{What exact divisibility already enforces}
\begin{proposition}[Consequences of an integral tail]
\label{prop:auditconsequences}
Suppose $D=MB-A>0$ and $D\mid A-1$, and put $q=1+(A-1)/D$.
The integer $q$ need not be prime. Nevertheless
\begin{align}
 &\gcd(D,A)=\gcd(D,MB)=1,\label{eq:auditcoprimeD}\\
 &\gcd(A,MB)=1,\label{eq:auditcoprimeA}\\
 &\gcd(q,A)=\gcd(D-1,A),\label{eq:auditgcdtail}\\
 &v_\ell(A-1)=v_\ell(D)+v_\ell(q-1)
 \quad\text{for every prime }\ell.\label{eq:auditvaluations}
\end{align}
If the full prime and ordering tests of
\eqref{eq:auditcertificate} also hold, then $D<B/2$, $D$ is odd and
$D\ge3$. If $M=2$, then $D\ge5$ and $3\nmid D$.
\end{proposition}
\begin{proof}
A common divisor of $D$ and $A$ divides both $A-1$ and $A$ and is
therefore $1$. Since $MB=A+D$, this also proves
$\gcd(D,MB)=1$ and $\gcd(A,MB)=1$. Reducing $Dq=MB-1$ modulo
$A$ gives $Dq\equiv D-1\pmod A$. Multiplication by the invertible
residue class $D$ preserves the greatest common divisor with $A$;
this proves \eqref{eq:auditgcdtail}. Equation
\eqref{eq:auditvaluations} is the valuation of the exact product
$A-1=D(q-1)$, with multiplicities included.

For a prime ordered tail, Lemma~\ref{lem:auditrounding} gives
$D<B/2$. Since $MB$ is even and $A$ is odd, $D$ is odd. If $D=1$,
then $q=A$. When $S$ has one element this violates $q>r$; when it
has more than one element it violates primality. Thus $D\ge3$.

More generally, if $3\mid D$, then \eqref{eq:auditcoprimeD} and
\eqref{eq:auditcoprimeA} show that no $p\in S$ is $0$ or $1$
modulo $3$. All are consequently $2$ modulo $3$. Writing $s=|S|$,
$D\mid A-1$ implies $(-1)^s\equiv1\pmod3$, so $s$ is even.
It follows that $B\equiv1\pmod3$. The equation $MB=A+D$ then
gives $M\equiv1\pmod3$. For $M=2$ this is impossible. Together
with oddness and the exclusion of $D=1$, it gives $D\ge5$.
\end{proof}

The equalities \eqref{eq:auditvaluations} are consequences of the
single identity $A-1=D(q-1)$, and \eqref{eq:auditcoprimeA} contains
the condition that no prime of $S$ divides another prime of $S$ minus
one. They prune a search, but they say nothing about the primality of
$q$.

\subsection{The deletion obstruction}
\begin{proposition}[No largest-prime deletion descent]
\label{prop:auditdeletion}
Suppose $n=Aq$ is a Lehmer composite and $q$ is its largest prime.
With $M=(n-1)/\varphi(n)$ and $B=\varphi(A)$, one has
\begin{equation}\label{eq:auditdelete}
 M-1<\frac{A-1}{B}<M.
\end{equation}
In particular $\varphi(A)$ does not divide $A-1$.
\end{proposition}
\begin{proof}
Set $D=MB-A$. By Lemma~\ref{lem:auditrounding}, $0<D<B/2$.
The integer $D$ is odd and $B$ is even. Thus $B>2$ and
$D\le B/2-1$, whence
\[
 0<\frac{D+1}{B}\le\frac12<1.
\]
Since $A=MB-D$, exact division gives
\[
 \frac{A-1}{B}=M-\frac{D+1}{B}.
\]
The last two displays give \eqref{eq:auditdelete}. There is no
integer strictly between the consecutive integers $M-1$ and $M$,
so the quotient cannot be an integer.
\end{proof}

So deleting the largest prime of a solution never gives a smaller
solution, whatever the quotient. A descent of another kind would have
to specify its new set of primes and prove both that they are prime and
that the totient condition holds; the minimality of a counterexample
does not provide either.

\subsection{Two remaining primes: factorisation and its exact inverse}
For this subsection $A,B$ describe an earlier prefix; $M\ge2$ is
fixed and $C=MB-A>0$. Consider distinct primes $r<p<q$ satisfying
\begin{equation}\label{eq:audittwotarget}
 Apq-1=MB(p-1)(q-1).
\end{equation}
Expanding the right side and moving terms gives, successively,
\begin{align*}
 Cpq-MB(p+q)+MB+1&=0,\\
 C^2pq-CMB(p+q)+C(MB+1)&=0,\\
 (Cp-MB)(Cq-MB)&=(MB)^2-C(MB+1).
\end{align*}
Using $MB-C=A$, the last right side is
\begin{equation}\label{eq:auditK}
 K=MAB-C=A^2+C(A-1).
\end{equation}
It is positive for $A\ge1$. For a genuine extension define
\begin{equation}\label{eq:audituv}
 u=Cp-MB,\qquad v=Cq-MB.
\end{equation}
The one-tail equation after choosing $p$ gives
$uq=MB(p-1)-1>0$, so $u>0$; then $v>u$ because $q>p$.
We have obtained
\begin{equation}\label{eq:auditfactorpair}
 uv=K,\qquad 0<u<v,\qquad u\equiv v\equiv-MB\pmod C.
\end{equation}

Conversely, any pair of positive integers satisfying
\eqref{eq:auditfactorpair} recovers integers
\begin{equation}\label{eq:auditrecover}
 p=\frac{u+MB}{C},\qquad q=\frac{v+MB}{C}.
\end{equation}
Reversing the displayed expansions proves
\eqref{eq:audittwotarget}. It is an extension of the prefix exactly
when $p$ and $q$ are primes greater than $r$.

\begin{figure}[tbp]
\centering
\begin{tikzpicture}[>=Latex,scale=0.95,every node/.style={font=\small}]
 \draw[->] (0,0)--(6.5,0) node[right] {$u/\sqrt K$};
 \draw[->] (0,0)--(0,4.3) node[above] {$v/\sqrt K$};
 \draw[blue!70!black,thick,domain=0.45:3.8,samples=100]
 plot ({1.5*\x},{1.5/\x});
 \draw[dashed,gray] (0,0)--(3.8,3.8);
 \fill[blue!70!black] (1.5,1.5) circle (2pt);
 \node[above right] at (1.5,1.5) {$u=v=\sqrt K$};
 \node[align=left,text width=5.1cm] at (8.2,2.8)
 {$uv=K$ is an exact factorisation.\\[3pt]
 Congruences give integral recovery.\\[3pt]
 Ordering and primality remain separate requirements.};
 \node[below] at (1.5,0) {$1$};
 \node[left] at (0,1.5) {$1$};
\end{tikzpicture}
\caption{The hyperbola $uv=K$, normalised. The two last primes come
from divisor pairs $(u,v)$ in the class \eqref{eq:auditfactorpair}
that recover primes.}
\label{fig:audithyperbola}
\end{figure}

\subsection{The extreme divisor}
For positive real $u\le\sqrt K$, the function $u+K/u$ decreases as
$u$ increases: its derivative is $1-K/u^2\le0$. Moreover
\begin{equation}\label{eq:auditproduct}
 pq=\frac{K+MB(u+K/u)+(MB)^2}{C^2}.
\end{equation}
So among the divisor pairs, the one with the larger smaller divisor
gives the smaller product $pq$.

The extreme divisor $u=1$ illustrates the distinction exactly. If
$C\mid MB+1$, then $C\mid A+1$. Write $j=(A+1)/C$, so that the
recovered first entry is $p=j+1$. Formula \eqref{eq:audituv}, or
\eqref{eq:audittwotarget}, yields
\[
 u(q-1)=Ap-1.
\]
With $u=1$ it follows that $q=Ap$. For $A>1$ this $q$ is composite,
whether or not $p$ is prime, so the divisor $u=1$ never gives a
solution. Other divisor pairs can give integer completions: the prefix $A=15$, $B=8$, $M=2$
gives $C=1$, $K=239$ and $(p,q)=(17,255)$. The exact equality
\[
 15\cdot17\cdot255-1=2\cdot8\cdot16\cdot254=65024
\]
holds, but $255=3\cdot5\cdot17$ is composite, so this is a spoof
solution in the sense of Section~\ref{sec:open} and not a solution.

\subsection{The outstanding universal assertion}
\begin{problem}[Universal prime-tail exclusion]\label{prob:universalprime}
Prove that for every nonempty finite set $S$ of distinct odd primes,
the three tests in \eqref{eq:auditcertificate} cannot all hold.
\end{problem}
By Theorem~\ref{thm:auditcriterion} this assertion is equivalent to
Lehmer's conjecture. The theorems of the paper prove it for the sets
$S$ with fewer than $15$ elements, with fewer than $20000$ elements
when $\lceil A/B\rceil\ge3$, and with fewer than $10^9$ elements when
$3\in S$, and for the sets with $A\cdot(1+(A-1)/D)\le10^{35}$.
